\chapter{Presentation de l'entreprise et cadre du stage}

\section{Presentation generale de l'entreprise d'accueil}

Le stage a ete realise au sein de \EntrepriseAccueil. Dans la version finale de ce rapport, cette section doit presenter de maniere synthetique l'entreprise, son secteur d'activite, sa taille, son positionnement et son organisation generale. Lorsque certaines informations relevent de la confidentialite, il est possible de conserver une presentation sobre, centree sur la nature des activites et sur le contexte technologique utile a la comprehension du stage.

\begin{longtable}{>{\raggedright\arraybackslash}p{4cm} >{\raggedright\arraybackslash}p{8.8cm}}
\toprule
Element & Information a renseigner \\
\midrule
Nom de l'entreprise & \EntrepriseAccueil \\
Secteur d'activite & [A completer] \\
Type de structure & [A completer] \\
Localisation & [A completer] \\
Service d'accueil & [A completer] \\
Nature du projet & Plateforme intelligente de gestion des tickets support \\
\bottomrule
\end{longtable}

Pour une redaction academique complete, cette partie peut couvrir :

\begin{itemize}[leftmargin=1.2cm]
  \item le domaine d'activite principal de l'entreprise ;
  \item son type de clientele ou ses utilisateurs internes ;
  \item son organisation generale ou son departement systemes d'information ;
  \item les enjeux de qualite, de support ou de digitalisation dans lesquels s'inscrit le projet.
\end{itemize}

Dans le cas de SupportFlow, l'environnement d'accueil est pertinent car il confronte directement l'etudiant a un besoin transversal : mieux piloter les interactions entre utilisateurs, agents support, managers et outils techniques. Le projet ne nait donc pas d'une logique purement scolaire. Il prend sens dans un besoin d'industrialisation du support et de clarification des responsabilites.

\subsection{Conseil de redaction}

Pour finaliser cette section, il est recommande de partir d'un texte court en trois temps :

\begin{enumerate}[leftmargin=1.2cm]
  \item presenter l'entreprise et son activite principale ;
  \item presenter le service ou l'equipe d'accueil ;
  \item expliquer le lien direct entre cette equipe et le besoin SupportFlow.
\end{enumerate}

\section{Equipe d'accueil et positionnement du projet}

Le projet s'inscrit generalement dans une equipe technique ou fonctionnelle liee a l'exploitation applicative, a la qualite de service, au support ou a la transformation digitale. Cette equipe constitue le point de contact entre les utilisateurs, les outils internes, les administrateurs et les responsables metier.

SupportFlow se positionne a l'intersection de plusieurs besoins :

\begin{itemize}[leftmargin=1.2cm]
  \item centraliser les demandes de support ;
  \item mieux qualifier les incidents et les demandes ;
  \item donner une visibilite claire au management ;
  \item rendre le traitement plus tracable et mesurable ;
  \item capitaliser les documents, rapports et resolutions ;
  \item preparer une trajectoire d'automatisation et d'assistance intelligente.
\end{itemize}

Ce positionnement explique la richesse du perimetre fonctionnel retenu. Le projet ne traite pas uniquement la saisie des tickets. Il couvre aussi le workflow, la securite, la priorisation, l'escalade, l'archivage, la notification et la qualite logicielle.

\subsection{Synthese a personnaliser}

\begin{quote}
[A adapter] Le projet SupportFlow a ete realise dans un contexte ou la gestion des demandes de support devait gagner en tracabilite, en rapidite et en capacite de supervision. L'equipe d'accueil avait besoin d'un outil capable de centraliser les tickets, de clarifier les responsabilites, de structurer les delais de traitement et de produire des preuves de suivi exploitables.
\end{quote}

\section{Mission confiee}

La mission confiee pendant le stage peut etre formulee comme suit : \textit{concevoir, realiser, corriger et consolider une plateforme intelligente de support capable de piloter le cycle de vie des tickets de maniere securisee, tracable et demonstrable.}

Cette mission se declinait en plusieurs sous-objectifs :

\begin{itemize}[leftmargin=1.2cm]
  \item analyser le besoin metier et identifier les acteurs du systeme ;
  \item definir l'architecture applicative et les technologies adequates ;
  \item developper les couches frontend et backend ;
  \item integrer Camunda, Keycloak, Alfresco, SonarQube et la chaine GitOps ;
  \item mettre en place des parcours coherents pour le client, l'agent, le manager et l'administrateur ;
  \item assurer la validation du produit a travers des tests, des scenarios et des preuves techniques.
\end{itemize}

\subsection{Formulation academique conseillee}

Dans la version finale du rapport, la mission peut etre reformulee en un paragraphe continu de huit a douze lignes, en repondant a trois questions :

\begin{itemize}[leftmargin=1.2cm]
  \item quel probleme concret devait etre traite ;
  \item quelle solution etait attendue ;
  \item quelle a ete votre contribution personnelle dans cette solution.
\end{itemize}

\section{Role de l'etudiant}

Le role de l'etudiant dans ce stage a ete large et structurant. Il ne s'est pas limite a l'execution d'une tache ponctuelle. Il a combine plusieurs responsabilites complementaires :

\begin{itemize}[leftmargin=1.2cm]
  \item \textbf{role d'analyste}, pour reformuler le besoin, decouper les parcours et identifier les regles metier ;
  \item \textbf{role de concepteur}, pour choisir l'architecture, les composants et les integrations ;
  \item \textbf{role de developpeur full stack}, pour realiser le frontend Angular, le backend Spring Boot et les integrations associees ;
  \item \textbf{role d'integrateur}, pour faire converger Camunda, Keycloak, Alfresco, l'agent IA, Docker, SonarQube et GitHub Actions ;
  \item \textbf{role de testeur et de consolideur}, pour corriger les anomalies, stabiliser les parcours, verifier les permissions et fiabiliser la demonstration ;
  \item \textbf{role de redacteur technique}, pour documenter le produit, preparer la soutenance et structurer les preuves de fonctionnement.
\end{itemize}

Cette polyvalence constitue l'un des apports majeurs du stage. Elle montre que le travail accompli ne porte pas uniquement sur du code, mais aussi sur la coherence globale d'un systeme metier.

\subsection{Point de vigilance redactionnel}

Dans un rapport de stage, il est important de distinguer clairement :

\begin{itemize}[leftmargin=1.2cm]
  \item ce qui relevait de la mission globale du projet ;
  \item ce qui relevait de votre travail personnel ;
  \item ce qui relevait d'outils, de bibliotheques ou de briques deja existantes.
\end{itemize}

Cette distinction renforce la credibilite du rapport et aide le jury a comprendre la vraie valeur de votre contribution.

\section{Cadre temporel et organisation du stage}

Le stage s'est deroule sur la periode \PeriodeStage. La conduite du projet a necessite une organisation rigoureuse, car plusieurs chantiers ont avance en parallele : realisation des fonctionnalites, correction des incoherences, gestion de la demonstration, stabilisation des environnements, mise en place de la qualite et preparation des livrables.

Le travail a suivi une logique progressive :

\begin{enumerate}[leftmargin=1.2cm]
  \item comprehension du besoin et cadrage du sujet ;
  \item construction du socle technique et fonctionnel ;
  \item enrichissement metier et integrations avancees ;
  \item corrections transverses et mise en coherence des parcours ;
  \item validation, documentation et preparation de la soutenance.
\end{enumerate}

\section{Enjeux du stage}

Le stage presentait plusieurs enjeux. Le premier etait \textbf{metier} : il fallait produire une plateforme utile et lisible pour plusieurs profils. Le deuxieme etait \textbf{technique} : il fallait faire cooperer de nombreuses briques heterogenes sans perdre la robustesse de l'ensemble. Le troisieme etait \textbf{pedagogique} : il fallait montrer une maitrise reelle des concepts, des choix d'architecture et des validations. Enfin, le quatrieme etait \textbf{demonstratif} : le produit devait pouvoir etre presente de facon claire, avec des preuves tangibles de qualite, de securite et de deploiement.

En cela, SupportFlow constitue un sujet de stage complet. Il met l'etudiant dans une situation proche d'un projet professionnel reel, ou la qualite du resultat depend autant de la technique que de la capacite a clarifier le metier, stabiliser les integrations et documenter les choix.

\section{Remarque sur la personnalisation finale}

Cette version du rapport a volontairement conserve des placeholders afin de permettre une personnalisation finale rapide sans remettre en cause la structure academique du document. Avant depot final, il conviendra de remplacer les champs suivants :

\begin{itemize}[leftmargin=1.2cm]
  \item nom complet de l'etudiant ;
  \item nom de l'etablissement et de la specialite ;
  \item nom de l'entreprise d'accueil ;
  \item periode exacte du stage ;
  \item nom des encadrants academique et entreprise ;
  \item donnees non confidentielles permettant de presenter correctement le contexte d'accueil.
\end{itemize}
